
\exo

Dans chaque cas, tracer dans un repère la droite dont on donne un point et un vecteur directeur.
\vspace{-6pt}
\setlength{\columnsep}{1cm}
\setlength{\columnseprule}{0pt}
\begin{multicols}{2}
\begin{enumerate}
\item $A\couple{4}{-1}$ et $\vecu\couple{3}{1}$.
\smallskip
\item $A\couple{0}{3}$ et $\vecu\couple{2}{-3}$.

\item $A\couple{-2}{0}$ et $\vecu\couple{0}{3}$.
\smallskip
\item $A\couple{3}{1}$ et $\vecu\couple{3}{0}$.
\end{enumerate}
\end{multicols}
\vspace{-6pt}

\exo

Dans chaque cas, déterminer une équation cartésienne de la droite $\cD$.
\begin{enumerate}
\item $\cD$ est la droite passant par $A\couple{-2}{1}$ et de vecteur directeur $\vecu\couple{3}{2}$.
\smallskip
\item $\cD$ est la droite $(AB)$, avec $A\couple{1}{-3}$ et $B\couple{-2}{1}$.
\smallskip
\item $\cD$ est la droite $(AB)$, avec $A\couple{0}{-5}$ et $B\couple{-3}{2}$.
\smallskip
\item $\cD$ est la droite passant par l'origine du repère et de vecteur directeur $\vecu\couple{3}{-2}$.
\smallskip
\item $\cD$ est la droite $(AB)$, avec $A\couple{2}{3}$ et $B\couple{-1}{3}$.
\smallskip
\item $\cD$ est la droite $(AB)$, avec $A\couple{1}{3}$ et $B\couple{1}{-2}$.
\end{enumerate}

\exo

Dans chaque cas, tracer dans un repère la droite $\cD$ dont on donne une équation cartésienne.
\vspace{-6pt}
\setlength{\columnsep}{1cm}
\setlength{\columnseprule}{0pt}
\begin{multicols}{3}
\begin{enumerate}
\item $\cD$ : $-3x+2y+1=0$.
\medskip
\item $\cD$ : $4x+5y-2=0$.
\medskip
\item $\cD$ : $\frac{3}{4}x-2y+5=0$.
\medskip
\item $\cD$ : $x-2=0$.
\medskip
\item $\cD$ : $2x-5y=0$.
\medskip
\item $\cD$ : $x-y=0$.
\medskip
\item $\cD$ : $x-3y=-3$.
\medskip
\item $\cD$ : $2y-3x+4=0$.
\medskip
\item $\cD$ : $-4x+y=1$.
\end{enumerate}
\end{multicols}
\vspace{-6pt}

\exo

Soient les points $A\couple{-2}{-3}$ et $B\couple{4}{-1}$.
\begin{enumerate}
\item Déterminer une équation cartésienne de la droite $(AB)$.
\item Le point $C\couple{3}{-1}$ appartient-il à la droite $(AB)$ ?
\item Déterminer l'ordonnée du point $D$, d'abscisse $\frac{3}{2}$, appartenant à $(AB)$.
\item Déterminer l'abscisse du point $E$, d'ordonnée $-\frac{4}{3}$, appartenant à $(AB)$.
\end{enumerate}

\exo

Soit $\cD$ la droite d'équation cartésienne $\frac{3}{4}x+y-\frac{1}{3}=0$.

Déterminer les coordonnées des points $A$ et $B$, intersections de $\cD$ avec les axes du repère.

\exo

\vspace{-12pt}
\setlength{\columnsep}{1cm}
\setlength{\columnseprule}{0pt}
\begin{multicols}{2}
\raggedcolumns
Soient un point $A\couple{x_A}{y_A}$ et une droite $\cD$ d'équation cartésienne $ax+by+c=0$.

On a commencé à rédiger en Python une fonction \pythoninline{test}, qui indique \pythoninline{"Oui"} si le point $A$ appartient à la droite $\cD$, et \pythoninline{"Non"} dans le cas contraire.

Compléter cette fonction.

\begin{pythoncodenumligne}
def test(xA, yA, a, b, c):
    if ........................:
        return "Oui"
    else:
        return "Non"
\end{pythoncodenumligne}
\end{multicols}
\vspace{-6pt}

\exo

Dans chaque cas, déterminer deux vecteurs directeurs de la droite dont on donne une équation.
\vspace{-6pt}
\setlength{\columnsep}{1cm}
\setlength{\columnseprule}{0pt}
\begin{multicols}{2}
\begin{enumerate}
\item $\cD$ a pour équation $-x+4y-7=0$.
\smallskip
\item $\cD$ a pour équation $3x-2y=0$.
\smallskip
\item $\cD$ a pour équation $4x-4y+5=0$.
\smallskip
\item $\cD$ a pour équation $5y+1=0$.
\smallskip
\item $\cD$ a pour équation $4x-5=0$.
\smallskip
\item $\cD$ passe par $A\couple{3}{3}$ et $B\couple{1}{7}$.
\smallskip
\item $\cD$ passe par $A\couple{2}{0}$ et est parallèle à $(Oy)$.
\smallskip
\item $\cD$ passe par $A\couple{2}{0}$ et est parallèle à $(Ox)$.
\end{enumerate}
\end{multicols}
\vspace{-6pt}

\exo

\vspace{-12pt}
\setlength{\columnsep}{1cm}
\setlength{\columnseprule}{0pt}
\begin{multicols}{2}
\raggedcolumns
On considère les six droites tracées ci-contre, et les six équations données ci-dessous :
\begin{enumerate}
\item $x-4=0$.
\item $-2x+y+3=0$.
\item $-x-3y=0$.
\item $y+2=0$.
\item $x+2y-3=0$.
\item $x-y+2=0$.
\end{enumerate}
Associer à chaque droite l'équation qui lui correspond.
\begin{center}
\def\xmin{-4}
\def\xmax{6}
\def\ymin{-3}
\def\ymax{4}
\def\xunite{0.85cm}
\def\yunite{0.85cm}
\psset{unit=\xunite, xunit=\xunite, yunit=\yunite, algebraic=true, dimen=middle, dotstyle=*, dotsize=3pt 0, linewidth=0.8pt, arrowsize=3pt 2, arrowinset=0.25, comma=true}
\begin{pspicture*}(\xmin,\ymin)(\xmax,\ymax)
\psgrid[xunit=\xunite, yunit=\yunite, subgriddiv=0, gridlabels=0, gridcolor=lightgray](0,0)(\xmin,\ymin)(\xmax,\ymax)
\psaxes[labelFontSize=\scriptstyle, xAxis=true, yAxis=true, Dx=1, Dy=1, ticksize=-2pt 0, subticks=1]{->}(0,0)(-3.99,-2.99)(\xmax,\ymax)
\psplot[plotpoints=200]{\xmin}{\xmax}{x+2}
\psplot[plotpoints=200]{\xmin}{\xmax}{-2}
\psplot[plotpoints=200]{\xmin}{\xmax}{1.5-0.5*x}
\psplot[plotpoints=200]{\xmin}{\xmax}{-x/3}
\psline(4,\ymin)(4,\ymax)
\psplot[plotpoints=200]{\xmin}{\xmax}{2*x-3}
\begin{footnotesize}
\uput{4pt}[135](1,3){$\cD_1$}
\uput{4pt}[270](2,-2){$\cD_2$}
\uput{4pt}[060](-3,3){$\cD_3$}
\uput{4pt}[060](-3,1){$\cD_4$}
\uput{4pt}[000](4,2){$\cD_5$}
\uput{4pt}[120](3,3){$\cD_6$}
\end{footnotesize}
\end{pspicture*}
\end{center}
\end{multicols}
\vspace{-12pt}

\exo

Dans chaque cas, tracer dans un repère la droite $\cD$ dont on donne l'équation réduite.
\vspace{-6pt}
\setlength{\columnsep}{1cm}
\setlength{\columnseprule}{0pt}
\begin{multicols}{3}
\begin{enumerate}
\item $\cD$ : $y=-5x+4$.
\medskip
\item $\cD$ : $y=\frac{3}{4}x-2$.
\medskip
\item $\cD$ : $x=-2$.
\medskip
\item $\cD$ : $y=\frac{7}{5}x+1$.
\medskip
\item $\cD$ : $y=3$.
\medskip
\item $\cD$ : $y=-x$.
\end{enumerate}
\end{multicols}
\vspace{-6pt}

\exo

Dans chaque cas, déterminer l'équation réduite de la droite $\cD$.
\begin{enumerate}
\item $\cD$ est la droite passant par $A\couple{2}{3}$ et de coefficient directeur $-3$.
\smallskip
\item $\cD$ est la droite passant par $A\couple{-2}{4}$ et de coefficient directeur $\frac{4}{5}$.
\smallskip
\item $\cD$ est la droite $(AB)$, avec $A\couple{-3}{-1}$ et $B\couple{5}{-3}$.
\smallskip
\item $\cD$ est la droite $(AB)$, avec $A\couple{-2}{1}$ et $B\couple{-2}{4}$.
\end{enumerate}

\exo

\vspace{-12pt}
\setlength{\columnsep}{1cm}
\setlength{\columnseprule}{0pt}
\begin{multicols}{2}
\raggedcolumns
Pour chacune des droites représentées ci-contre, déterminer par lecture graphique son coefficient directeur et son ordonnée à l'origine, puis donner son équation réduite.
\begin{center}
\def\xmin{-4}
\def\xmax{6}
\def\ymin{-3}
\def\ymax{4}
\def\xunite{0.85cm}
\def\yunite{0.85cm}
\psset{unit=\xunite, xunit=\xunite, yunit=\yunite, algebraic=true, dimen=middle, dotstyle=*, dotsize=3pt 0, linewidth=0.8pt, arrowsize=3pt 2, arrowinset=0.25, comma=true}
\begin{pspicture*}(\xmin,\ymin)(\xmax,\ymax)
\psgrid[xunit=\xunite, yunit=\yunite, subgriddiv=0, gridlabels=0, gridcolor=lightgray](0,0)(\xmin,\ymin)(\xmax,\ymax)
\psaxes[labelFontSize=\scriptstyle, xAxis=true, yAxis=true, Dx=1, Dy=1, ticksize=-2pt 0, subticks=1]{->}(0,0)(-3.99,-2.99)(\xmax,\ymax)
\psplot[plotpoints=200]{\xmin}{\xmax}{1-x/3}
\psplot[plotpoints=200]{\xmin}{\xmax}{1.5*x-2}
\psplot[plotpoints=200]{\xmin}{\xmax}{3-2*x}
\psplot[plotpoints=200]{\xmin}{\xmax}{-2}
\begin{footnotesize}
\uput{4pt}[080](-3,2){$\cD_1$}
\uput{4pt}[120](3,2.5){$\cD_2$}
\uput{4pt}[225](-0.5,4){$\cD_3$}
\uput{4pt}[080](-3,-2){$\cD_4$}
\end{footnotesize}
\end{pspicture*}
\end{center}
\end{multicols}
\vspace{-12pt}

\exo

Dans chaque ligne du tableau ci-dessous, $A$ et $B$ sont deux points d'une droite $\cD$, $\vecu$ est un vecteur directeur de $\cD$ et $m$ est le coefficient directeur de $\cD$.

Lorsque c'est possible, compléter ce tableau.
\begin{center}
\setlength{\tabcolsep}{0.1cm}
\renewcommand{\arraystretch}{1.5}
\begin{tabular}{
*{4}{|>{\centering\arraybackslash}m{1.8cm}}
|>{\centering\arraybackslash}m{3.2cm}|}
\hline 
Point $A$ & Point $B$ & $\vecu$ & $m$ & Équation de $\cD$ \\
\hline
$\couple{3}{9}$ & $\couple{2}{4}$ & & & \\
\hline 
 &  & & & $y=4$ \\
\hline 
 & $\couple{3}{1}$ & & $\frac{5}{3}$ & \\
\hline 
$\couple{3}{1}$ & & $\couple{7}{3}$ & & \\
\hline 
\end{tabular} 
\end{center}

